% Generated draft: requires mathematical and source review.
\documentclass[11pt]{article}
\usepackage[margin=25mm]{geometry}
\usepackage{fontspec}
\setmainfont{Latin Modern Roman}
\usepackage{amsmath,amssymb,mathtools,unicode-math}
\setmathfont{Latin Modern Math}
\usepackage{enumitem,needspace,xurl,hyperref}
\setlength{\parindent}{0pt}
\setlength{\parskip}{5pt}
\setlength{\emergencystretch}{3em}
\newcommand{\N}{\mathbb{N}}
\newcommand{\Z}{\mathbb{Z}}
\newcommand{\Q}{\mathbb{Q}}
\newcommand{\R}{\mathbb{R}}
\newcommand{\C}{\mathbb{C}}
\newcommand{\F}{\mathbb{F}}
\newcommand{\A}{\mathbb{A}}
\newcommand{\PP}{\mathbb{P}}
\newcommand{\E}{\mathbb{E}}
\newcommand{\eps}{\varepsilon}
\newcommand{\dd}{\,\mathrm{d}}
\begin{document}
\section*{UnsolvedMath ID: 4800023}
\textbf{Source problem code:} catalogue entry
\subsection{Definitions and mathematical statement}
A Hardy field is a field of germs at $+\infty$ of real-valued smooth functions, closed under differentiation. Let $\ell\geq1$ and let $a_1,\ldots,a_\ell$ belong to one Hardy field and have polynomial growth: for each $i$ there exist $C_i>0$, $k_i\geq0$ with $|a_i(t)|\leq C_i t^{k_i}$ for all sufficiently large $t$.

Interpret a nonzero linear combination as one with a nonzero coefficient vector. Thus assume that for every $\boldsymbol\beta=(\beta_1,\ldots,\beta_\ell)\in\mathbb R^\ell\setminus\{0\}$, every $c\in\mathbb R$, and every $p\in\mathbb Z[t]$,
\[
\lim_{t\to\infty}\frac{\left|\sum_{i=1}^\ell\beta_i a_i(t)-cp(t)\right|}{\log t}=+\infty.
\]
In particular, this convention excludes a nontrivial coefficient relation giving the zero germ. If the source's $\operatorname{span}^*(\mathcal F)$ instead removes the zero function after taking the span, that is a different hypothesis and needs clarification.

For every ergodic invertible probability-preserving system $(X,\mathcal B,\mu,T)$ and every $f_1,\ldots,f_\ell\in L^\infty(\mu)$, prove
\[
\lim_{N\to\infty}\left\|
\frac1N\sum_{n=1}^N\prod_{i=1}^\ell T^{\lfloor a_i(n)\rfloor}f_i
-\prod_{i=1}^\ell\int_X f_i\,d\mu
\right\|_{L^2(\mu)}=0.
\]
Here $T^k f=f\circ T^k$, $\lfloor u\rfloor$ is the greatest integer at most $u$, and the product of integrals is viewed as a constant function. Invertibility is needed for negative exponents, which the stated growth and separation hypotheses do not exclude. The supplied research summary attributes this Hardy-field joint-convergence conclusion to Bergelson--Moreira--Richter (2024).
\subsection{Short English statement}
If every nontrivial real combination of polynomial-growth Hardy-field functions stays farther than a logarithmic scale from every real multiple of an integer polynomial, their integer-part multiple ergodic averages should converge in $L^2$ to the product of the integrals. The supplied record reports a 2024 resolution.
\subsection{Sources}
\begin{itemize}
\item Nikos Frantzikinakis, \emph{Some open problems on multiple ergodic averages} (2016 source-list version), Problem 23, \url{https://arxiv.org/abs/1103.3808}; accessed 2026-07-29, with source-TeX extraction recorded in the supplied provenance.
\item Bergelson--Moreira--Richter, \emph{Advances in Mathematics} 443 (2024), article 109597, as the resolution citation supplied in the research summary. No article title is given in the record.
\item The background refers to the author's November 2025 progress page as not recording a resolution of this numbered problem; no URL for that page is supplied.
\item Local registry \texttt{data/unsolvedmath/problems.json}, permanent UnsolvedMath ID $4800023$, source code catalogue entry, \emph{Multiple ergodic averages --- Problem 23}. Dataset attribution: Ulam AI and the UnsolvedMath Contributors.
\end{itemize}
\end{document}
